\documentclass[UTF8,fontset=none,11pt,a4paper]{ctexart}
\usepackage{ipara-handbook}
\handbooksetup{DS-A01 序列扫描与区间状态}{408 · 数据结构算法 Extension}{Prefix · Window · Pointer · Boundary}
\begin{document}
\handbooktitle{序列扫描与区间状态}{怎样让边界移动、摘要复用和候选淘汰都有证明}{408 · 数据结构 Algorithm Extension A01}

\begin{corebox}{母模型：已确认区、候选区与永久淘汰}
序列技巧的共同链条是
\[
\begin{aligned}
\text{Order / Coordinate Contract}
&\to\text{Active Interval or Result Prefix}
\to\text{Incremental Summary}\\
&\to\text{Dominance / Boundary Proof}
\to\text{One-Pass Cost}.
\end{aligned}
\]
指针移动只有在被跳过的元素以后不可能重新成为答案时才安全；摘要只有在能从旧摘要和新端点更新时才值得维护。模板名称不是正确性证明。
\end{corebox}

\begin{methodbox}{Owner 与停止边界}
\begin{itemize}
\item \kw{Owns：}前缀/差分、窗口、快慢/左右指针、矩阵边界和候选支配。
\item \kw{Uses：}DS01 顺序/矩阵表示、DS02 单调队列、DS09 有序索引、DS10 Hash 计数。
\item \kw{Stops：}二分答案的单调谓词归 DS09；链表指针改边归 DS01；单调栈/队列容器不变量归 DS02。
\end{itemize}
\end{methodbox}
\tableofcontents\newpage

\section{前缀和与差分：把重复区间工作变成端点事件}

对数组 $a_0,\ldots,a_{n-1}$，定义半开前缀
\[
P[0]=0,\qquad P[i+1]=P[i]+a_i.
\]
闭区间 $[l,r]$ 的和为 $P[r+1]-P[l]$。$n+1$ 个位置的哨兵使 $l=0$、空区间和边界统一，不需要为第一项写特殊分支。二维前缀把一维差分扩展为容斥：
\[
P[i+1][j+1]=a_{ij}+P[i][j+1]+P[i+1][j]-P[i][j].
\]
查询矩形时以四个前缀加减，左上角被减两次所以加回。

差分是前缀的逆操作：
\[
D[0]=a_0,\qquad D[i]=a_i-a_{i-1}.
\]
对闭区间 $[l,r]$ 加 $v$ 只改两个端点：$D[l]+=v$，若 $r+1<n$ 则 $D[r+1]-=v$；最后做一次前缀恢复。$m$ 次修改加恢复的成本是 $O(m+n)$，但修改过程中不能直接读取最终数组，除非另有在线结构。

\begin{examplebox}{半开事件避免端点混乱}
把乘客在 $from$ 上车、在 $to$ 下车表示为事件 $+num$ 于 $from$、$-num$ 于 $to$，扫描站点时先累加事件再检查容量。这里区间是 $[from,to)$；若题意规定在终点之后才离开，则应改成 $to+1$。端点语义属于问题合同，不能死记某一道题的下标。
\end{examplebox}

前缀技巧依赖所用聚合可逆：加法可以用相减恢复，异或可用再次异或恢复；一般的 `min`/`max` 没有这样的逆运算，不能机械套两个前缀相减。溢出、负数、空数组和二维矩阵的行列方向都必须进入边界检查。

\subsection{前缀值频次：把“区间等式”改写成“历史查询”}

若目标是统计和为 $k$ 的连续子数组，令扫描到当前位置后的前缀为 $s$。任何以当前位置结尾的合法区间都对应一个历史前缀 $s-k$，所以维护
\[
\code{freq}[x]=\text{此前前缀值 }x\text{ 出现的次数}.
\]
处理当前元素时必须先把 \code{freq[s-k]} 加入答案，再把当前 $s$ 写入频次；反过来会在 $k=0$ 时把“当前前缀减自身”的空区间错误计入。初始化 \code{freq[0]=1} 表示数组开始前的空前缀，使从下标 $0$ 开始的合法区间无需特判。该方法允许负数，因为它不依赖窗口和的单调性，期望时间 $O(n)$、空间 $O(n)$，其中 Hash 的期望成本前提归 DS10。

同一个改写还生成两类变体。若要找最长区间，Hash 保存某前缀值的\kw{最早位置}而不是出现次数，且第一次出现后不覆盖；若二元数组要求 $0$ 与 $1$ 数量相等，可先把 $0$ 映射为 $-1$，问题变成“和为零的最长区间”。统计数量、求最长长度和判断存在性使用的是不同摘要，不能只因都叫“前缀 + Hash”就共享更新合同。

\subsection{左右摘要并不要求可逆：除自身以外的乘积}

有些题不做任意区间查询，而是为每个位置组合“左侧全部信息”和“右侧全部信息”。除自身以外数组乘积可令输出先保存左侧乘积，再从右向左用一个滚动后缀乘积相乘，时间 $O(n)$、除输出外空间 $O(1)$，且无需除法，因此自然覆盖零元素。它不是两个前缀相除：乘法遇到零不可逆。这里可复用的是\kw{两侧摘要分解}，不是前缀和公式。

\subsection{差分的密集数组与稀疏事件}

当坐标范围小而连续，可用长度固定的差分数组承载事件。全零初始数组无需先做相邻差；航班编号从 $1$ 开始时，闭区间 $[first,last]$ 应映射为 $D[first-1]{+}{=}v$、若 $last<n$ 则 $D[last]{-}{=}v$。拼车则是半开行程 $[from,to)$，在 $to$ 直接减。两个公式的差别来自业务端点，不来自题名。

当坐标很大但事件很少，不应为最大坐标分配整段数组；可把“坐标 $\to$ 净变化”存入有序事件表，按坐标排序后扫描。会议室并发数、日程最大重叠等问题都可由事件前缀生成，但同一坐标的开始/结束先后必须由半开或闭区间合同决定。若任务要求在线区间修改与在线查询，离线差分已经越过停止边界，应转向 Fenwick/线段树等动态结构，而不是在每次查询前都 $O(n)$ 恢复。

\section{滑动窗口：可恢复的活动区间}

窗口维护 $[left,right)$ 和一个随端点增删的摘要。典型框架是右端加入元素，若窗口违反约束则左端不断移出，窗口重新合法后更新答案。正确性依赖一个单调性条件：在右端固定时，若某个左端已使窗口合法/非法，向某一方向移动左端不会出现不可预测的多次来回。若约束不满足该条件，双指针可能漏解。

\subsection{最长合法窗口与最短覆盖窗口}

无重复最长子串维护字符计数；右端加入导致计数超过 1 时，左端移除直到恢复合法，每个字符至多进出一次，所以总时间 $O(n)$。最小覆盖窗口的状态不只是计数，还要维护“已满足的种类数”或缺口数；只有当所有需求满足时才收缩并记录最短答案。固定长度排列窗口则在每次移动时增删一个字符，不能把“满足后继续收缩”误用于固定 $k$。

可变窗口的完整状态协议是 \code{need/window/valid}：\code{need[c]} 表示目标需求，\code{window[c]} 表示当前计数，只有当某字符计数\kw{首次达到}需求时才增加 \code{valid}，从达到状态跌破时才减少。最小覆盖在窗口已经覆盖全部需求时，于移除左端\kw{之前}记录候选；无重复最长子串在恢复合法后记录最大长度。更新时机若互换，前者会错过即将被破坏的最短合法窗，后者会把仍含重复的窗口算入答案。

固定长度排列/异位词窗口保持 $right-left=|T|$：右端每加入一项，就在超长时恰好移出左端，再判断计数是否完全匹配。它不需要寻找“最短”，也不允许在合法后任意继续缩小。窗口最大值虽然长度也固定，但摘要是 DS02 Own 的单调队列；A01 只负责索引过期条件 $index<left$ 和端点推进。

\begin{warnbox}{什么时候不能用普通窗口}
“和为 $k$ 的最长子数组”在有负数时不具备按右端扩张、左端收缩的单调性；窗口会错过解，应改用前缀和 + Hash 或其他状态。窗口模板适合的是可由端点单调修复的约束，而不是所有子数组题。
\end{warnbox}

滑动窗口最大值需要 DS02 的单调队列：窗口端点负责过期，队尾弹出被新值支配的索引，队头才是最大值。队列实现若使用线性 `shift()`，代码层面的总复杂度会被破坏，应使用头指针或真正的双端队列。

\section{快慢指针：结果前缀与扫描后缀}

原地筛选数组时，`fast` 扫描未处理后缀，`slow` 指向结果区的下一个写入位置。循环不变量是：$[0,slow)$ 已经是按题意保留的结果，$[fast,n)$ 尚未读取；若当前值满足保留条件，就写入 `nums[slow]` 并推进，否则只推进 `fast`。因此删除元素、去重和移动零都能在 $O(n)$、$O(1)$ 辅助空间完成。

覆盖与交换不是同一语义：覆盖只承诺前缀有效，尾部可以是垃圾；移动零若要求保留非零元素相对顺序，可以用交换或稳定覆盖后补零，必须在合同中声明。

若有序数组允许每个值最多保留 $q$ 次，写指针为 $slow$ 时，只需在 $slow<q$ 或当前值不等于结果前缀中倒数第 $q$ 个值时写入：
\[
\code{keep}(x)\iff slow<q\lor x\ne a[slow-q].
\]
这个式子把“去重一次”和“最多保留两次”统一为结果前缀约束；比较对象必须来自已经确认的输出区，而不是未处理原数组。含退格字符串的逆向扫描则维护待跳过字符数：遇退格增加计数，遇普通字符且计数非零就消耗一次跳过；两串各找下一个有效字符后比较。它仍是“扫描后缀逐步提交结果”的边界模型，但写入方向和状态不同。

链表快慢指针属于 DS01 的具体表示应用。本册只抽取“速度比导致相遇/中点”的推理：快指针每次两步、慢指针一步，若有环则相对速度使它们在有限圈内相遇；相遇后一个指针回到头部、两个同速前进，利用距离等式得到环入口。空链、单节点和 `n` 超长时的先后指针必须由 dummy 与边界条件覆盖。

\section{左右指针：单调性与支配淘汰}

左右指针从两端向内收缩，只有在某一侧移动后能永久排除一批候选时才正确。

\subsection{有序两数和与三数和}

有序数组中若 $a_l+a_r<target$，固定 $r$ 时任何更小的左值都不可能达到目标，因此可以安全地增加 $l$；若和过大则对称减少 $r$。三数和先排序并固定一个元素，把剩余问题降为左右指针；外层和内层都必须跳过相邻重复值，避免输出重复而不误删不同位置的合法组合。

$k$ 数之和可以递归固定一个值，把 $k$ 降为 $k-1$，直到两数和；每层只跳过\kw{同一递归深度}的相邻重复候选，命中一组两数和后同时越过左右两侧的重复值。若把“值相等”不分层级地全局跳过，会删掉需要多个相同数且索引不同的合法答案。最接近目标的三数和则不在命中后收集所有组合，而是维护当前绝对误差最小值，并仍用和的单调方向移动端点。排序成本 $O(n\log n)$，三数主体 $O(n^2)$；固定更高 $k$ 后成本按递归层数增长，输出规模还需单独报告。

\subsection{容器与接雨水}

容器面积为 $(r-l)\min(h_l,h_r)$。宽度每步都减少，若左边较短，移动右边不能提高当前瓶颈高度且宽度变小，故右端在这一轮被支配；只能移动短边寻找更高瓶颈。接雨水维护左右最大高度，若 `leftMax < rightMax`，左侧水量已由 `leftMax` 决定，右侧未知不会改变它，因此结算左端并右移；等号时采用固定的一侧即可，但要保持不重计。

\section{矩阵：坐标变换与四边界收缩}

顺时针旋转正方形可分解为主对角线转置再逐行翻转；两步都是原地交换，组合后每个坐标 $(i,j)$ 映射到 $(j,n-1-i)$。螺旋遍历维护 `top,bottom,left,right` 四个边界，每走完上、右、下、左边就收缩对应边界；走下边和左边前必须再次检查边界，才能避免单行/单列重复。

由坐标映射可复原三种旋转，而不必分别背代码：
\[
\begin{array}{c|c|c}
\text{变换}&(i,j)\text{ 的新位置}&\text{一种原地分解}\\\hline
90^\circ\text{ 顺时针}&(j,n-1-i)&\text{主对角线转置 + 逐行翻转}\\
90^\circ\text{ 逆时针}&(n-1-j,i)&\text{主对角线转置 + 逐列翻转}\\
180^\circ&(n-1-i,n-1-j)&\text{逐行翻转 + 逐列翻转}
\end{array}
\]
这些原地分解要求方阵；一般 $m\times n$ 矩阵旋转后形状变成 $n\times m$，若存储不能重解释形状，就需要新容器。

螺旋读取与螺旋生成共享同一 frontier：前者把边界元素追加到结果，后者把递增值写入边界。每条边完成后立即收缩，因此在只剩一行或一列时，下边/左边必须重新检查合法性。对角线遍历则用 $i-j$ 标识主对角线平行族、用 $i+j$ 标识副对角线平行族；若输出方向交替，还需由对角线编号奇偶决定反转，而不是改变分组不变量。

矩阵边界算法的循环不变量是：边界外已输出/已写入，边界内尚未处理；每一轮至少收缩一条边，故每个格子最多访问一次。对角线编号由 $i-j$ 或 $i+j$ 不变量给出；矩阵置零、生命游戏等原地题还需要额外编码“旧状态”和“新状态”，不能把已经修改的值误当作邻居旧值。

\begin{methodbox}{矩阵原地修改的两阶段提交}
若新状态依赖周围旧状态，扫描阶段必须同时保留 old/new 两层语义：可以使用额外位、枚举过渡态或第一行/列作为标记。只有全部单元都完成读取后，第二遍才把过渡态提交为最终状态。直接边读边覆盖会让后续邻居读到混合时间层。
\end{methodbox}

\section{题型路由：先识别可复用信息，再选模板}

\begin{center}\small
\begin{tabularx}{\linewidth}{L{3.2cm}YY}
\toprule
题面工作负载 & 首选状态 & 停止/转交信号\\\midrule
静态多次区间聚合 & 前缀摘要 & 聚合不可逆或需要在线修改\\
离线多次区间增量 & 端点差分事件 & 稀疏坐标改事件表；在线查询改动态树\\
连续区间且约束可单调修复 & 滑动窗口 & 负数等使收缩方向失去单调性\\
原地稳定筛选/压缩 & 结果前缀 + 扫描后缀 & 尾部也有语义时补写/交换合同\\
有序候选的配对/支配 & 左右指针 & 无序且不能排序，或移动侧无法证明淘汰\\
二维边界访问/坐标置换 & 四边界或坐标映射 & 新状态依赖旧邻居时采用两阶段提交\\
单调可行答案空间 & DS09 二分谓词 & 判定函数不单调\\
\bottomrule
\end{tabularx}\end{center}

\section{复杂度与反例验证}

\begin{center}\small
\begin{tabularx}{\linewidth}{L{3cm}YY}
\toprule
模式 & 典型成本 & 必须证明的事\\\midrule
前缀/差分 & 构建/恢复 $O(n)$，查询或修改 $O(1)$ & 哨兵和端点语义\\
窗口 & 双端点总移动 $O(n)$ & 收缩方向不会漏解\\
快慢筛选 & $O(n)$、$O(1)$ 辅助空间 & 结果前缀不变量\\
左右指针 & 常见 $O(n)$ 或排序后 $O(n^2)$ & 移动侧的支配淘汰\\
矩阵边界 & $O(mn)$ & 边界外完成、单行单列不重复\\
\bottomrule
\end{tabularx}\end{center}

验证时至少攻击：全相等、严格递增、严格递减、空/单元素、重复端点、负数窗口、奇偶长度、单行/单列矩阵和超过整数范围。若改变一个前提后模板失效，应把前提写进 Stop Boundary，而不是继续套用。

\section{压缩复原}

忘记代码时问：当前已确认区是什么？活动边界采用闭区间还是半开区间？摘要加入/移出如何更新？被丢弃元素凭什么永不成为答案？窗口约束是否单调？矩阵每圈收缩哪条边？答案必须能复原循环不变量、端点公式和复杂度。

\begin{corebox}{一句话压缩}
序列算法的核心不是“两根指针”，而是维护一个可解释的边界与摘要，并证明每次移动都扩大已确认区域或永久淘汰候选；没有单调性、可逆性或支配关系时，应停止套模板并换用 DP、Hash、排序或二分等 Owner。
\end{corebox}
\end{document}
